\section{How to Read the Numbers}

The mathematics in this paper is elementary enough to check by hand. The difficulty is not the arithmetic; it is keeping several kinds of arithmetic from being blended together. This section supplies the small set of tools needed to follow the main findings.

\subsection{Gematria Begins with the Hebrew Form}

Standard gematria assigns numerical values to the Hebrew letters: aleph through tet carry 1 through 9, yod through tsadi carry the tens from 10 through 90, and qof through tav carry 100 through 400. The value of a word or phrase is the sum of its letters.\cite{schechter1906gematria}

For example, YHWH is written yod-he-vav-he:

\[
10+5+6+5=26.
\]

The spelling is part of the data. An article, a vav, a final letter, or a different word division can change the result. For that reason, this study values the displayed consonantal form rather than a silently normalized version of it.

\subsection{Mispar Gadol Adds a Second Value Layer}

Mispar Gadol gives expanded values to the five letters that have special final forms:

\begin{center}
\begin{tabular}{lr@{\qquad}lr@{\qquad}lr}
\toprule
Final letter & Value & Final letter & Value & Final letter & Value \\
\midrule
kaf & 500 & mem & 600 & nun & 700 \\
pe & 800 & tsadi & 900 & & \\
\bottomrule
\end{tabular}
\end{center}

This is not a replacement for standard gematria. It is a second view of the same form, used only when a final letter is actually present. The term \emph{Mispar Gadol} has been used for more than one large-value procedure; throughout this paper it means the convention in which final forms continue the hundreds series from 500 through 900.\cite{schechter1906gematria} Some patterns appear in both systems; Elohim, for example, has standard value 86 but Mispar Gadol value 646, which is a decimal palindrome.

\subsection{A Base Is a Way of Writing a Number}

Changing bases does not change the underlying number. It changes its notation. In base 10, place values are powers of 10. In base 12, they are powers of 12. Thus

\[
22_{12}=2(12)+2=26_{10}.
\]

This is why YHWH can be written as both $26_{10}$ and $22_{12}$. The equality is numerical; the visible 22 is a property of the base-12 notation.

Base labels matter whenever two strings are compared. Decimal 22 and $22_{12}$ look alike but are different values:

\[
22_{10}=1\mathrm{A}_{12},
\qquad
22_{12}=26_{10}.
\]

I use A for the base-12 digit ten and B for eleven. A result such as $7\mathrm{A}7_{12}$ is therefore a three-digit base-12 number, not a mixture of letters and decimal digits.

\subsection{Digit Sums: The 8 Layer and the Exact 12 Layer}

The original study repeatedly summed digits until the result was 12 or less. Under that declared rule, for example,

\[
386\rightarrow3+8+6=17\rightarrow1+7=8.
\]

The 8 endpoint is robust: changing the stopping threshold from 9 through 15 leaves the same 10 primary rows touching 8. The language around 12 needs a distinction. A conventional digital root continues until one digit remains, so 12 cannot be a conventional final root. In this corpus, however, eight primary rows produce 12 immediately on the first stated digit sum. That exact event remains true regardless of what one chooses to do afterward.

I therefore use two precise descriptions:

\begin{itemize}
\item \textbf{terminal 8}: the declared repeated sum reaches 8;
\item \textbf{exact first-step 12-sum}: the first sum of the relevant displayed digits equals 12.
\end{itemize}

For a base-12 expression, the displayed base-12 digits are summed as values, with A equal to 10 and B equal to 11. Yeshua's $282_{12}$ is an exact example because $2+8+2=12$.

\subsection{Palindromes Depend on Their Base}

A palindrome reads the same forward and backward. Since a palindrome is a property of a written string, changing the base can create or remove one. The value 314 is not a decimal palindrome, but

\[
314_{10}=222_{12},
\]

so Shaddai becomes palindromic in base 12. The main primary examples are:

\begin{itemize}
\item YHWH: $22_{12}$;
\item Adonai: $55_{12}$;
\item Shaddai: $222_{12}$;
\item Ehyeh Asher Ehyeh: $393_{12}$;
\item El Olam under Mispar Gadol: $515_{12}$.
\end{itemize}

Yeshua's $282_{12}$ and Yeshua HaMashiach's $525_{12}$ belong to declared Christian layers. HaKadosh Baruch Hu contributes $7\mathrm{A}7_{12}$ in a later-traditional layer. The scroll-construction value $6556_{12}$ is also palindromic, but it is a constructed canonical object rather than a literal name.

I treat palindrome symmetry as an exact observation. What that symmetry represents, whether identity, stability, beauty, selection, or chance, is an interpretive question.

\subsection{Reference Values and Mathematical Echoes}

Two kinds of numerical reference appear in the study.

\paragraph{Canonical and internal references.}
The values 22, 42, and 72 belong to the religious and textual world of the paper. \emph{Sefer Yetzirah} calls the Hebrew letters the ``Twenty-Two Letters of Foundation''; \emph{Kiddushin} 71a names a forty-two-letter divine Name; and the later seventy-two-name construction is drawn from the three 72-letter verses of Exodus 14:19--21.\cite{seferyetzirah22,kiddushin42,bahya72} The scroll construction supplies a separate internal value, $6556_{12}$.

\paragraph{Ordinary mathematical echoes.}
Some values resemble familiar mathematical objects. Shaddai's 314 can be read as the digits of $3.14$, while El Shaddai's 345 presents the ordered digits of the 3-4-5 Pythagorean triple. Other star-number, Mersenne-form, and Fibonacci-adjacent observations appear in the full derivations. These are genuine digit or value relationships, but their relevance to divine names is less constrained than the canonical references. I include them as invitations for further thought, not as load-bearing evidence.

\subsection{Intersections Rather Than Boxes}

The central idea of the paper can now be stated simply. A name may carry several independent features at once:

\begin{itemize}
\item a terminal 8;
\item an exact first-step sum of 12;
\item a decimal palindrome;
\item a base-12 palindrome;
\item an exact shared value or canonical digit-string contact.
\end{itemize}

The most interesting rows are not necessarily those that fit one category most cleanly. They are the rows where several of these features meet.
